% Code-aligned method draft for the current spatial_combination implementation.

\documentclass[10pt]{article}

\usepackage[margin=1in]{geometry}
\usepackage[T1]{fontenc}
\usepackage[utf8]{inputenc}
\usepackage{amsmath,amssymb,mathtools}
\usepackage{bm}
\usepackage{booktabs}
\usepackage{graphicx}
\usepackage{microtype}
\usepackage[hidelinks]{hyperref}

\newcommand{\R}{\mathbb{R}}
\newcommand{\one}{\mathbb{1}}
\newcommand{\cL}{\mathcal{L}}
\newcommand{\cT}{\mathcal{T}}
\newcommand{\cC}{\mathcal{C}}
\newcommand{\cP}{\mathcal{P}}
\newcommand{\cS}{\mathcal{S}}
\newcommand{\eps}{\varepsilon}
\newcommand{\sg}{\operatorname{sg}}
\newcommand{\Pool}{\operatorname{Pool}}
\newcommand{\Norm}{\operatorname{Norm}}
\newcommand{\clip}{\operatorname{clip}}
\newcommand{\softplus}{\operatorname{softplus}}
\newcommand{\stopgrad}{\operatorname{stopgrad}}
\newcommand{\BCE}{\operatorname{BCE}}
\newcommand{\Dice}{\operatorname{Dice}}
\newcommand{\KL}{\operatorname{KL}}
\newcommand{\sigm}{\operatorname{sigmoid}}

\title{OODKA-OT}
\date{}

\begin{document}
\maketitle

\section{Problem Setting and Notation}

% Code anchors: oodka/data/slice_dataset.py; oodka/models/feature_extraction.py;
%               oodka/train/model_builder.py

Let $B_{\phi_b}$ be a frozen BiomedParse foundation model and $E_{\phi_n}$
be a frozen nnUNet expert.  The subscripts $b$ and $n$ below refer to
BiomedParse and nnUNet, respectively.  The trainable parameters are collected in
$\Theta$ and contain only the feature adapters, two-branch decomposition modules,
and the spatial prompt router; $\phi_b$ and $\phi_n$ remain fixed.  We write
$C_{\mathrm{pr}}$ for the number of text prompts, reserving $P$ and $S$ for the
two feature branches.

Training uses batches of $B$ independent blocks, each containing $Z$ contiguous
center slices.  The nnUNet tensor is
$\bm{x}^{n}\in\R^{B\times C_n\times Z\times H\times W}$, whereas the
BiomedParse tensor is
$\bm{x}^{b}\in\R^{B\times Z\times 3\times H\times W}$.  A 2D nnUNet and the
2D BiomedParse backbone both process the $BZ$ slices as a flattened batch.  Their
features are then restored to five-dimensional tensors.  At the four implemented
levels
\begin{equation}
    \mathcal{I}=\{2,3,4,5\},
\end{equation}
we denote the foundation and expert features by
\begin{equation}
    \bm{Z}^{b}_{\ell}\in
    \R^{B\times C^b_\ell\times Z\times H_\ell\times W_\ell},
    \qquad
    \bar{\bm{Z}}^{n}_{\ell}\in
    \R^{B\times C^n_\ell\times Z\times \bar H_\ell\times \bar W_\ell}.
\end{equation}
Before channel adaptation, define
\begin{equation}
    \bm{Z}^{n}_{\ell}
    =\operatorname{Resize}_{\mathrm{tri}}
      (\bar{\bm{Z}}^{n}_{\ell};Z,H_\ell,W_\ell),
\end{equation}
where resizing is the identity when the shapes already agree.  This resized
tensor is both the expert-adapter input and its reconstruction target.  Repeated
slices in an incomplete tail block are marked by
$v_{bz}=0$, assigned target label $-1$, and excluded from the segmentation,
reconstruction, separation, and OT terms.  The route KL is defined once per
block and prompt field and is therefore independent of the number of valid
slices.

For generic 2.5D training, the three BiomedParse channels of center slice $z$ are
$(x_{z-1},x_z,x_{z+1})$, with boundary replication.  The current LGE pathway uses
$Z=1$ and repeats the center slice, $(x_z,x_z,x_z)$.

For the generic CT path, the default BiomedParse intensity transform uses window
level $40$ and window width $400$.  Equivalently,
\begin{equation}
    x_i^b=255\,
    \frac{\clip(x_i,-160,240)+160}{400}.
\end{equation}
Both the CT transform and the LGE transform in Eq.~\eqref{eq:mri-normalization}
produce values in $[0,255]$ before the frozen BiomedParse pixel normalization.

\section{Multi-scale Two-Branch Decomposition}

% Code anchors: oodka/models/disentangle.py;
%               oodka/train/forward.py::_disentangle_and_inject

We use $P$ and $S$ to name the two learned branches.  Their operational roles
are determined by the objectives below: $P$ receives structure-anchored balanced
OT supervision, and $S$ receives task-aware residual UOT supervision.

\subsection{Foundation-side decomposition}

At each level, two independent bias-free $1\!\times\!1\!\times\!1$
convolutions decompose the frozen BiomedParse feature:
\begin{equation}
    \widehat{\bm{Z}}^{b}_{\ell,p}
        =D^{b}_{\ell,p}(\bm{Z}^{b}_{\ell}),
    \qquad
    \widehat{\bm{Z}}^{b}_{\ell,s}
        =D^{b}_{\ell,s}(\bm{Z}^{b}_{\ell}).
    \label{eq:base-decomposition}
\end{equation}
Both outputs preserve the channel count and spatial shape of
$\bm{Z}^{b}_{\ell}$.  No normalization or activation is applied after these
projections.

\subsection{Expert-side channel adaptation and decomposition}

The resized nnUNet feature first passes through a trainable channel adapter
\begin{equation}
    \bm{U}_{\ell}
    =\operatorname{GELU}\!\left(
       \operatorname{IN}\!\left(A^0_\ell(\bm{Z}^{n}_{\ell})\right)
     \right),
\end{equation}
where all $A_\ell$ are bias-free $1\!\times\!1\!\times\!1$ convolutions.
The two expert components are
\begin{equation}
    \widehat{\bm{Z}}^{n}_{\ell,p}=N_{\ell,p}(A_{\ell,p}\bm{U}_{\ell}),
    \qquad
    \widehat{\bm{Z}}^{n}_{\ell,s}=N_{\ell,s}(A_{\ell,s}\bm{U}_{\ell}).
    \label{eq:expert-decomposition}
\end{equation}
Here $N_{\ell,p}$ and $N_{\ell,s}$ are affine instance normalization in the
default configuration, except that both are identities at level $\ell=5$.
Separate decoders $R_{\ell,p}$ and $R_{\ell,s}$, each implemented as
Conv--IN--GELU--Conv, map the components back to $C^n_\ell$ channels.

\subsection{Feature and decoder reconstruction}

For a target tensor $\bm{T}$ and a broadcastable valid-slice mask $\bm{V}$,
define the energy-normalized mean squared error
\begin{equation}
    \operatorname{NMSE}(\bm{A},\bm{T};\bm{V})
    =
    \frac{
      \sum_i V_i(A_i-T_i)^2/N_V
    }{
      \sum_i V_iT_i^2/N_V+10^{-8}
    },
    \qquad
    N_V=C H W\sum_{b,z}v_{bz}.
    \label{eq:nmse}
\end{equation}
The two reconstruction terms at level $\ell$ are
\begin{align}
    \cL^{n}_{\mathrm{rec},\ell}
    &={\operatorname{NMSE}}\!\left(
      R_{\ell,p}(\widehat{\bm{Z}}^{n}_{\ell,p})
      +R_{\ell,s}(\widehat{\bm{Z}}^{n}_{\ell,s}),
      \bm{Z}^{n}_{\ell};\bm{V}\right),\\
    \cL^{b}_{\mathrm{rec},\ell}
    &={\operatorname{NMSE}}\!\left(
      \widehat{\bm{Z}}^{b}_{\ell,p}
      +\widehat{\bm{Z}}^{b}_{\ell,s},
      \bm{Z}^{b}_{\ell};\bm{V}\right).
    \label{eq:level-reconstruction}
\end{align}

The frozen BiomedParse pixel decoder is also evaluated three times: once with
all $P$ features injected, once with all $S$ features injected, and once with
their level-wise sums injected.  Let
$\bm{f}^{r}_{p}$, $\bm{f}^{r}_{s}$, and $\bm{f}^{r}_{p+s}$ denote the resulting
features at $r\in\mathcal{Q}=\{\text{mask},5,4,3\}$.  The decoder consistency
term is
\begin{equation}
    \cL_{\mathrm{pd}}
    =\frac{1}{4}\sum_{r\in\mathcal{Q}}
      \operatorname{NMSE}
      (\bm{f}^{r}_{p}+\bm{f}^{r}_{s},\bm{f}^{r}_{p+s};\bm{V}).
\end{equation}
The reconstruction loss used by the implementation is therefore
\begin{equation}
    \cL_{\mathrm{ae}}
    =\frac{1}{9}\left[
       \sum_{\ell\in\mathcal{I}}
       \left(\cL^{n}_{\mathrm{rec},\ell}
             +\cL^{b}_{\mathrm{rec},\ell}\right)
       +\cL_{\mathrm{pd}}
     \right].
    \label{eq:ae-total}
\end{equation}
Thus, the eight level-wise terms are summed, the four pixel-decoder terms are
first averaged into one term, and the resulting nine terms are averaged.

\subsection{Cross-correlation separation}

For $\bm{P},\bm{S}\in\R^{B\times C\times Z\times H\times W}$, flatten the
valid spatial dimensions and standardize each channel to obtain
$\widetilde{\bm{P}},\widetilde{\bm{S}}\in\R^{B\times C\times N}$.  The
implemented separation penalty is the mean absolute channel cross-correlation
\begin{equation}
    \operatorname{CorrAbs}(\bm{P},\bm{S})
    =\frac{1}{BC^2}\sum_{b,c,c'}
       \left|
       \frac{1}{N_b}\sum_{i=1}^{N_b}
       \widetilde P_{bci}\widetilde S_{bc'i}
       \right|.
\end{equation}
When every slice is valid, standardization uses PyTorch's channel-wise sample
standard deviation.  When a tail mask is present, the implementation instead
uses the valid-site population RMS after masked mean centering and divides the
correlation by the valid-site count $N_b$.
It is applied to both models at all four levels and summed:
\begin{equation}
    \cL_{\mathrm{ort}}
    =\sum_{\ell\in\mathcal{I}}
      \left[
      \operatorname{CorrAbs}(\widehat{\bm{Z}}^{b}_{\ell,p},
                              \widehat{\bm{Z}}^{b}_{\ell,s})
      +
      \operatorname{CorrAbs}(\widehat{\bm{Z}}^{n}_{\ell,p},
                              \widehat{\bm{Z}}^{n}_{\ell,s})
      \right].
    \label{eq:ortho}
\end{equation}

\section{Prompt-conditioned Spatial Beta Re-fusion}

% Code anchors: oodka/models/beta_router.py;
%               oodka/models/biomedparse_helpers.py;
%               oodka/train/forward.py::_fuse_all_prompt_features

Let $\bm{t}_c\in\R^{d_t}$ be the frozen BiomedParse class embedding of prompt
$c$.  The router first computes
\begin{equation}
    \bm{h}_c
    =\operatorname{GELU}\!\left(
       W_h\operatorname{LN}(\bm{t}_c)+\bm{b}_h
     \right).
\end{equation}
Two linear heads map $\bm{h}_c$ to coefficient vectors
$\bm{a}^{\alpha}_c,\bm{a}^{\beta}_c\in\R^K$.  The $K=G^2$ centers
$\{\bm{\mu}_k\}_{k=1}^{K}$ lie on a $G\times G$ grid in $[-1,1]^2$.  At
normalized coordinate $\bm{u}$, the radial basis is
\begin{equation}
    \varphi_k(\bm{u})
    =\exp\!\left(
      -\frac{\|\bm{u}-\bm{\mu}_k\|_2^2}{2\sigma_{\mathrm{rbf}}^2}
      \right).
\end{equation}
By default $G=8$ and
$\sigma_{\mathrm{rbf}}=1.5\,[2/(G-1)]=3/(G-1)$.
The prompt-specific concentration fields are
\begin{align}
    \eta_c^{\alpha}(\bm{u})
      &=b_{\alpha}+\frac{1}{\sqrt K}
        \sum_{k=1}^{K}a^{\alpha}_{ck}\varphi_k(\bm{u}),\\
    \eta_c^{\beta}(\bm{u})
      &=b_{\beta}+\frac{1}{\sqrt K}
        \sum_{k=1}^{K}a^{\beta}_{ck}\varphi_k(\bm{u}),\\
    \alpha_c(\bm{u})&=1+\softplus(\eta_c^{\alpha}(\bm{u})),\qquad
    \beta_c(\bm{u})=1+\softplus(\eta_c^{\beta}(\bm{u})).
    \label{eq:beta-field}
\end{align}

During training, a reparameterized sample is drawn independently for every
block and prompt,
\begin{equation}
    g_{bc}(\bm{u})\sim
    \operatorname{Beta}(\alpha_c(\bm{u}),\beta_c(\bm{u})).
\end{equation}
The same sampled field is shared by all $Z$ slices in that block.  At validation
and deployment, the deterministic mean is used:
\begin{equation}
    \bar g_c(\bm{u})
    =\frac{\alpha_c(\bm{u})}
           {\alpha_c(\bm{u})+\beta_c(\bm{u})}.
    \label{eq:beta-mean}
\end{equation}

The default prior is parameterized by mean $m_0=0.7$ and concentration
$\kappa_0=10$, hence $(\alpha_0,\beta_0)=(7,3)$.  Both coefficient heads are
zero-initialized, and the scalar biases are initialized so that every prompt and
position initially follows this prior.  The router regularizer is
defined on the mask-feature grid $\Omega_{\mathrm{mask}}$:
\begin{equation}
    \cL_{\mathrm{route}}
    =\frac{1}{C_{\mathrm{pr}}|\Omega_{\mathrm{mask}}|}
      \sum_{c=1}^{C_{\mathrm{pr}}}
      \sum_{\bm{u}\in\Omega_{\mathrm{mask}}}
      \KL\!\left[
        \operatorname{Beta}(\alpha_c(\bm{u}),\beta_c(\bm{u}))
        \,\|\,
        \operatorname{Beta}(\alpha_0,\beta_0)
      \right].
    \label{eq:route-kl}
\end{equation}

One field is generated at the mask-feature resolution and area-resized to all
three predictor feature resolutions.  For $r\in\mathcal{Q}$, the fused feature
for visual slice $(b,z)$ and prompt $c$ is
\begin{equation}
    \widetilde{\bm{f}}^{r}_{bzc}(\bm{u})
    =g^{r}_{bc}(\bm{u})\bm{f}^{r}_{p,bz}(\bm{u})
     +\left[1-g^{r}_{bc}(\bm{u})\right]
       \bm{f}^{r}_{s,bz}(\bm{u}).
    \label{eq:spatial-fusion}
\end{equation}
The gate is scalar across channels but varies over prompt and spatial position.
The $BZC_{\mathrm{pr}}$ visual--prompt pairs are flattened in $(b,z,c)$ order
and evaluated in one call to the frozen BiomedParse predictor.  If the predictor
produces $Q_m$
mask queries, their mask logits are averaged without object-existence weighting:
\begin{equation}
    \ell_{bcz}(\bm{u})
      =\frac{1}{Q_m}\sum_{h=1}^{Q_m}\ell^{(h)}_{bcz}(\bm{u}).
    \label{eq:query-average}
\end{equation}
The logits are finally trilinearly resized to the configured $(Z,H,W)$ output
shape.

\section{Prompt-wise Segmentation Objective}

% Code anchor: oodka/train/forward.py::_compute_segmentation_loss_and_metrics

For prompt $c$ associated with class identifier $k_c$, define the binary target
$y_{bcv}=\one[Y_{bv}=k_c]$ and valid-voxel mask
$m_{bv}=\one[Y_{bv}\neq-1]v_{bz(v)}$.  The per-pair BCE and soft Dice loss are
\begin{align}
    \cL^{\mathrm{bce}}_{bc}
    &=\frac{\sum_v m_{bv}\,
       \BCE(\ell_{bcv},y_{bcv})}
       {\max(\sum_v m_{bv},1)},\\
    \cL^{\mathrm{dice}}_{bc}
    &=1-\frac{
       2\sum_v m_{bv}\sigm(\ell_{bcv})y_{bcv}+10^{-6}}
       {\sum_v m_{bv}\sigm(\ell_{bcv})+
        \sum_v m_{bv}y_{bcv}+10^{-6}}.
\end{align}
A target is treated as empty when its foreground occupies less than
$5\times10^{-4}$ of the valid volume:
\begin{equation}
    e_{bc}
    =\one\!\left[
       \sum_vm_{bv}y_{bcv}<5\times10^{-4}\sum_vm_{bv}
      \right].
\end{equation}
The corresponding loss and weight are
\begin{equation}
    \ell^{\mathrm{seg}}_{bc}
      =\cL^{\mathrm{bce}}_{bc}
       +(1-e_{bc})\cL^{\mathrm{dice}}_{bc},
    \qquad
    w_{bc}=r_{bc}\left(1-\tfrac12 e_{bc}\right),
    \label{eq:pair-seg}
\end{equation}
where $r_{bc}\in\{0,1\}$ is the optional prompt-validity indicator.  Generic
training uses
\begin{equation}
    \cL_{\mathrm{seg}}^{\mathrm{mean}}
    =\frac{\sum_{b,c}w_{bc}\ell^{\mathrm{seg}}_{bc}}
           {\max(\sum_{b,c}w_{bc},10^{-6})}.
    \label{eq:seg-mean}
\end{equation}
The LGE branches use the implemented prompt-sum reduction
\begin{equation}
    \cL_{\mathrm{seg}}^{\mathrm{LGE}}
    =\sum_c
      \frac{\sum_bw_{bc}\ell^{\mathrm{seg}}_{bc}}
           {\max(\sum_bw_{bc},10^{-6})}.
    \label{eq:seg-lge}
\end{equation}

\section{Component-aware OT Distillation}

% Code anchors: oodka/models/ot/{cost,mass,sinkhorn,transport,losses,objective}.py;
%               oodka/train/forward.py::_compute_detached_pixel_error_maps

OT is evaluated independently at every $\ell\in\mathcal{I}$.  Invalid repeated
tail slices are removed, leaving $M$ valid 2D feature maps.  Each spatial
dimension is adaptively average-pooled to at most $32$, never upsampled, and the
result is flattened into tokens.  Let
$\bm{x}^{p}_i,\bm{x}^{s}_i\in\R^{C_\ell}$ be foundation tokens and
$\bm{y}^{p}_j,\bm{y}^{s}_j\in\R^{C_\ell}$ be channel-aligned expert tokens.
The sample index is suppressed below: the implementation solves $M$ independent
transport problems in a batched tensor, rather than pooling tokens from different
slices into one plan.

\subsection{Structure mass for the $P$ branch}

For class $c$, let $\bm{M}_c=\one[\bm{Y}=k_c]$ on a valid slice.  Three pooled
structure maps are constructed:
\begin{align}
    \bm{O}_c&=\operatorname{AvgPool}(\bm{M}_c),\\
    \bm{B}_c&=\operatorname{AvgPool}\!\left(
       \clip(\operatorname{Dilate}_{3\times3}(\bm{M}_c)
            -\operatorname{Erode}_{3\times3}(\bm{M}_c),0,1)
       \right),\\
    \bm{R}_c&=\operatorname{MaxPool}(\bm{M}_c).
\end{align}
The class budget is zero for an absent class and inverse-square-root in its area
otherwise:
\begin{equation}
    \bar\pi_c=
    \begin{cases}
      |\bm{M}_c|^{-1/2},&|\bm{M}_c|>0,\\
      0,&|\bm{M}_c|=0,
    \end{cases}
    \qquad
    \bm{\pi}=\Norm(\bar{\bm{\pi}}),
\end{equation}
where $\Norm$ uses a uniform vector if all entries are zero.  In that all-empty
case, the structure maps below are themselves zero and the implementation falls
back to feature-energy masses.
The unnormalized structure score at token $i$ is
\begin{equation}
    h_i=\sum_c\pi_c
       \left(O_{c,i}+1.5B_{c,i}+0.2R_{c,i}\right).
    \label{eq:structure-score}
\end{equation}
For a pooled feature $\bm{F}$, define its mean-normalized token energy
\begin{equation}
    \mathcal{E}_i(\bm{F})
    =\frac{\sqrt{\|\bm{F}_i\|_2^2+\eps}}
           {\max\!\left(
            \operatorname{mean}_{i'}
            \sqrt{\|\bm{F}_{i'}\|_2^2+\eps},\eps\right)}.
\end{equation}
The foundation demand and expert supply candidates are
\begin{equation}
    q_i^b=h_i\left(1+0.1\mathcal{E}_i(\bm{x}^p)\right),
    \qquad
    q_j^n=h_j\left(1+0.1\mathcal{E}_j(\bm{y}^p)\right).
\end{equation}
If $\sum_i h_i=0$, each candidate is replaced by its corresponding energy.
Finally,
\begin{equation}
    \bm{a}^p=\Norm(\bm{q}^b),\qquad
    \bm{b}^p=\Norm(\bm{q}^n),\qquad
    \Norm(\bm{q})=\frac{\bm{q}}{\sum_iq_i},
    \label{eq:p-mass}
\end{equation}
with a uniform fallback if the candidate mass is numerically zero.

\subsection{Task-aware residual mass for the $S$ branch}

The $S$ energy and its specificity relative to both components are
computed from the stabilized magnitude
$\nu(\bm{z})=\sqrt{\|\bm{z}\|_2^2+\eps}$:
\begin{equation}
    e_i=\frac{\nu(\bm{s}_i)}
              {\max(\operatorname{mean}_{i'}\nu(\bm{s}_{i'}),\eps)},
    \qquad
    r_i=\frac{\nu(\bm{s}_i)}
              {\nu(\bm{p}_i)+\nu(\bm{s}_i)+\eps}.
    \label{eq:s-energy-specificity}
\end{equation}
These quantities are computed separately for foundation and expert tokens: on
the foundation side $(\bm{p}_i,\bm{s}_i)=(\bm{x}^p_i,\bm{x}^s_i)$, and on the
expert side $(\bm{p}_j,\bm{s}_j)=(\bm{y}^p_j,\bm{y}^s_j)$.

For each prompt, a detached pixel-wise BCE map is computed from the current
foundation logit and the expert logit.  In grouped-label LGE training, the expert
logit for group $G_c$ is
\begin{equation}
    \ell^n_c
    =\operatorname{logit}\!\left(
       \clip\!\left[
       \sum_{k\in G_c}\operatorname{softmax}(\bm{\ell}^n)_k,
       10^{-6},1-10^{-6}
       \right]\right).
    \label{eq:grouped-expert-logit}
\end{equation}
The per-prompt errors are averaged over valid prompts and pooled to the OT grid,
giving foundation difficulty $d_i$ and expert error $e_i^{\mathrm{err}}$.  Their
advantage is
\begin{equation}
    A_i=d_i-e_i^{\mathrm{err}}.
\end{equation}
The current default uses a dense, bounded smooth-advantage score:
\begin{align}
    \widehat d_i
      &=\clip\!\left(
        \frac{d_i}{\max(\operatorname{mean}_{i'}d_{i'},\eps)},0,3
        \right),\\
    s_A&=\max\!\left(\operatorname{mean}_i|A_i|,\eps\right),\\
    g_i&=\clip\!\left(
       2\sigm\!\left(\frac{A_i}{0.5s_A}\right),0,3
       \right).
    \label{eq:smooth-gain}
\end{align}
Thus $g_i=1$ is neutral, $g_i>1$ indicates lower expert error, and $g_i<1$
indicates higher expert error.  The residual candidates and normalized marginals
are
\begin{align}
    q_i^b&=e_i^b r_i^b(0.1+\widehat d_i),&
    \bm{a}^s&=\Norm(\bm{q}^b),\\
    q_j^n&=e_j^n r_j^n(0.1+g_j),&
    \bm{b}^s&=\Norm(\bm{q}^n).
    \label{eq:s-mass}
\end{align}

\subsection{Transport costs}

Only normalized, detached copies of the pooled features enter the cost.  The
stabilized copies are
$\bar{\bm{x}}_i=\bm{x}_i/\max(\|\bm{x}_i\|_2,10^{-6})$ and analogously for
$\bar{\bm{y}}_j$; a zero token therefore remains zero.  The feature and
coordinate costs are
\begin{align}
    C^{\mathrm{feat}}_{ij}
      &=\left[1-\langle\bar{\bm{x}}_i,
                             \bar{\bm{y}}_j\rangle\right]_+,\\
    C^{\mathrm{coord}}_{ij}
      &=\left[\|\bm{u}_i-\bm{u}_j\|_2-r_0\right]_+^2,
      \qquad \bm{u}_i,\bm{u}_j\in[-1,1]^2.
    \label{eq:feature-coordinate-cost}
\end{align}
For the $P$ branch, pooled GT class vectors are $\ell_1$-normalized across
classes using the same stabilized-normalization convention, and the semantic
cost is
\begin{equation}
    C^{\mathrm{sem}}_{ij}
      =\left[1-\langle\bm{m}_i,\bm{m}_j\rangle\right]_+.
\end{equation}
The implemented costs are
\begin{align}
    \bm{C}^p
      &=\bm{C}^{\mathrm{feat}}
        +0.25\bm{C}^{\mathrm{coord}}
        +0.25\bm{C}^{\mathrm{sem}},\\
    \bm{C}^s
      &=\bm{C}^{\mathrm{feat}}
        +0.25\bm{C}^{\mathrm{coord}},
    \qquad r_0=0.25.
    \label{eq:ot-costs}
\end{align}

\subsection{Balanced $P$ transport and unbalanced $S$ transport}

The $P$ plan is the entropically regularized balanced OT solution
\begin{equation}
    \bm{T}^{p}
    =\arg\min_{\bm{T}\geq0}
       \langle\bm{T},\bm{C}^{p}\rangle
       +\eps_p\sum_{ij}T_{ij}(\log T_{ij}-1),
    \quad
    \bm{T}\bm{1}=\bm{a}^{p},\quad
    \bm{T}^{\top}\bm{1}=\bm{b}^{p}.
    \label{eq:balanced-ot}
\end{equation}
The $S$ plan relaxes both marginals with generalized KL penalties:
\begin{equation}
\begin{split}
    \bm{T}^{s}=\arg\min_{\bm{T}\geq0}\;&
       \langle\bm{T},\bm{C}^{s}\rangle
       +\eps_s\sum_{ij}T_{ij}(\log T_{ij}-1)\\
       &+\rho_b\KL(\bm{T}\bm{1}\,\|\,\bm{a}^{s})
       +\rho_n\KL(\bm{T}^{\top}\bm{1}\,\|\,\bm{b}^{s}).
    \label{eq:unbalanced-ot}
\end{split}
\end{equation}
The default values are $\eps_p=\eps_s=0.1$, $\rho_b=1.0$, and $\rho_n=0.2$.
Both plans are computed in log space for 30 Sinkhorn iterations.  For UOT, the
corresponding relaxation exponents are
\begin{equation}
    \tau_b=\frac{\rho_b}{\rho_b+\eps_s},\qquad
    \tau_n=\frac{\rho_n}{\rho_n+\eps_s}.
\end{equation}

Given either plan, the expert tokens are projected to foundation positions using
the barycentric teacher
\begin{equation}
    \widetilde{\bm{y}}_i
    =\frac{\sum_jT_{ij}\bm{y}_j}
           {\max(\sum_jT_{ij},10^{-8})}.
    \label{eq:barycentric-teacher}
\end{equation}
The weighted cosine distillation operator is
\begin{equation}
    \mathcal{D}(\bm{x},\widetilde{\bm{y}};\bm{w})
    =\frac{\sum_iw_i
       \left[1-\cos(\bm{x}_i,\widetilde{\bm{y}}_i)\right]}
       {\max(\sum_iw_i,10^{-8})}.
\end{equation}
Consequently,
\begin{align}
    \cL_{P\text{-OT},\ell}
       &=\mathcal{D}(\bm{x}^{p},\widetilde{\bm{y}}^{p};\bm{a}^{p}),\\
    \cL_{S\text{-UOT},\ell}
       &=\mathcal{D}(\bm{x}^{s},\widetilde{\bm{y}}^{s};
                     \bm{T}^{s}\bm{1}),\\
    \cL_{P\text{-OT}}&=\frac{1}{4}\sum_{\ell\in\mathcal{I}}
       \cL_{P\text{-OT},\ell},\qquad
    \cL_{S\text{-UOT}}=\frac{1}{4}\sum_{\ell\in\mathcal{I}}
       \cL_{S\text{-UOT},\ell}.
    \label{eq:ot-distillation-losses}
\end{align}
If the total UOT mass received by the foundation is at most $10^{-6}$, the
corresponding $S$ loss is set to zero.
Because every balanced marginal sums to one, valid slices are equally weighted
in the batched $P$ loss.  In the $S$ loss, valid slices are instead weighted by
their total UOT-received mass.

All masses, errors, costs, Sinkhorn plans, and barycentric teachers are detached.
The distillation gradient therefore updates the foundation-side $P/S$ adapters
through the raw foundation tokens, but it does not update the frozen expert or
differentiate through the transport solver.  The trainable expert-side adapters
receive gradients from reconstruction and separation, but not from the detached
OT teacher path.

\section{Overall Training Objective and Optimization}

% Code anchors: oodka/train/forward.py::forward_one_batch;
%               oodka/train/engine.py

For generic training, the per-batch objective is
\begin{equation}
    \cL
    =\lambda_{\mathrm{seg}}\cL_{\mathrm{seg}}
     +\lambda_{\mathrm{ae}}\cL_{\mathrm{ae}}
     +\lambda_{\mathrm{ort}}\cL_{\mathrm{ort}}
     +\lambda_{\mathrm{route}}\cL_{\mathrm{route}}
     +\lambda_p\cL_{P\text{-OT}}
     +\lambda_s\cL_{S\text{-UOT}}.
    \label{eq:overall-objective}
\end{equation}
The generic entry point defaults to
\begin{equation}
    (\lambda_{\mathrm{seg}},\lambda_{\mathrm{ae}},
     \lambda_{\mathrm{ort}},\lambda_{\mathrm{route}},
     \lambda_p,\lambda_s)
    =(3.0,0.2,0.3,10^{-3},0.1,0.1).
\end{equation}

The route weight is linearly increased over the first five epochs:
\begin{equation}
    \lambda_{\mathrm{route}}(e)
    =10^{-3}\min\!\left(1,\frac{e}{5}\right).
    \label{eq:route-schedule}
\end{equation}
For either OT target weight $\lambda^\star$ with start epoch $e_0$ and warm-up
length $E_w=5$, the applied value at epoch $e$ is
\begin{equation}
    \lambda(e)=
    \begin{cases}
      0,&e<e_0,\\
      \lambda^\star\min\!\left(1,
      \dfrac{e-e_0+1}{E_w}\right),&e\geq e_0.
    \end{cases}
    \label{eq:ot-schedule}
\end{equation}
The defaults are $e_0=2$ for $P$-OT and $e_0=3$ for $S$-UOT.

All trainable modules are optimized jointly with AdamW, base learning rate
$10^{-4}$, weight decay $10^{-4}$, automatic mixed precision when enabled, and
global gradient-norm clipping at $5$.  The generic entry point uses a constant
learning rate by default.  The LGE entry point uses epoch-wise cosine decay to
$0.05$ of the initial rate and no learning-rate warm-up.  For an $E$-epoch LGE
run, the exact epoch-level schedule is
\begin{equation}
    \eta_e=\eta_0\left[
      r_{\min}+(1-r_{\min})
      \frac{1+\cos\!\left(\pi\frac{e-1}{E-1}\right)}{2}
      \right],
    \qquad r_{\min}=0.05.
    \label{eq:lge-cosine-lr}
\end{equation}

\section{LGE Spatial Combination with Predicted-ROI Refinement}

% Code anchors: oodka/models/prompts.py; oodka/data/lge_roi.py;
%               oodka/train/lge_roi_engine.py; run_train_lge_roi.py;
%               run_eval_lge_roi.py

The current LGE pathway reuses the same BiomedParse backbone, decomposition
modules, pixel decoder, predictor, and Beta router in two forwards.  There are no
separate student parameters for the full-image and ROI branches.

\subsection{Aligned LGE preprocessing and label groups}

BiomedParse inputs are precomputed with the crop and resampling geometry from
the frozen nnUNet plans.  Only the intensity normalization is replaced.  Let
$q_{0.1}$ be the $0.1$th percentile of the volume and
$\Omega_{\mathrm{fg}}=\{i:x_i>q_{0.1}\}$.  The lower and upper values are the
first and $99$th percentiles over $\Omega_{\mathrm{fg}}$, and the normalized
image is
\begin{equation}
    x_i^{b}=255\,
      \frac{\clip(x_i,q_1,q_{99})-q_1}{q_{99}-q_1},
    \label{eq:mri-normalization}
\end{equation}
with $q_{99}$ replaced by $q_1+1$ in the degenerate case.  The result is clipped
to $[0,255]$ and stored as unsigned 8-bit data.

The source labels are scar $1$, edema $2$, LV cavity $3$, normal myocardium $4$,
and RV cavity $5$.  Pass~A uses the three disjoint anatomy groups
\begin{equation}
    \mathcal{G}^{A}=\bigl((3),(5),(1,2,4)\bigr),
    \label{eq:anatomy-groups}
\end{equation}
corresponding to LV, RV, and total myocardium.  The split-pathology refinement
implemented by the current V3 launcher uses
\begin{equation}
    \mathcal{G}^{R}_{\mathrm{split}}
    =\bigl((3),(5),(4),(1),(2)\bigr),
    \label{eq:split-groups}
\end{equation}
corresponding to LV, RV, normal myocardium, scar, and edema.  The V2 merged
variant replaces the last two groups by $(1,2)$:
\begin{equation}
    \mathcal{G}^{R}_{\mathrm{merged}}
    =\bigl((3),(5),(4),(1,2)\bigr).
    \label{eq:merged-groups}
\end{equation}
Grouping is applied independently to the targets of each pass; source groups are
required to be disjoint.  We denote the anatomy prompt count by $C_A=3$ and the
refinement prompt count by $C_R=5$ for the split variant or $C_R=4$ for the
merged variant.

\subsection{Pass A: full-image anatomy and ROI construction}

Pass~A predicts logits
$\bm{\ell}^{A}=(\ell^{A}_{\mathrm{LV}},\ell^{A}_{\mathrm{RV}},
\ell^{A}_{\mathrm{MYO}})$ on the full image.  The localization probability and
threshold set are
\begin{equation}
    p_{\mathrm{MYO}}(\bm{u})=\sigm(\ell^{A}_{\mathrm{MYO}}(\bm{u})),
    \qquad
    \Omega_{\tau}=\{\bm{u}:p_{\mathrm{MYO}}(\bm{u})\geq\tau\},
    \quad \tau=0.3.
    \label{eq:roi-mask}
\end{equation}
For nonempty $\Omega_\tau$, the smallest axis-aligned bounding box is computed.
Each side length is expanded by factor $\gamma=1.25$ about its center, with a
minimum side length of $8$ pixels, then shifted/clipped to the image boundary.
If $\Omega_\tau$ is empty, the default fallback is the full image.

The LGE model is first trained for $10$ epochs using only Pass~A.  Before epoch
$11$, deterministic student predictions generate one ROI per training slice.
With the default refresh interval of zero, this cache remains fixed for all
subsequent training epochs.  Validation and inference instead compute an online
ROI from the current Pass~A prediction.  ROI generation is discrete and detached;
no gradient propagates through thresholding, bounding-box extraction, or crop
coordinates.

\subsection{Pass R: crop, resize, and local refinement}

For every slice, the same ROI is applied to the nnUNet input, the BiomedParse
input, and the target.  Image crops are bilinearly resized and label crops are
nearest-neighbor resized back to the configured $256\times256$ model input.
Pass~R then predicts either four merged-pathology logits or five split-pathology
logits using the same student parameters as Pass~A.

For V2/V3 training, an ROI prompt-validity mask prevents a truncated positive
class from being interpreted as a reliable negative.  If source group $G_c$ is
absent from the complete slice, its prompt remains valid as a genuine negative.
If it is present, its refinement loss is enabled only when
\begin{equation}
    \frac{|\{Y\in G_c\}\cap\mathrm{ROI}|}
         {|\{Y\in G_c\}|}\geq0.01.
    \label{eq:roi-visibility}
\end{equation}
This indicator is the $r_{bc}$ used in Eq.~\eqref{eq:pair-seg} and in the
prompt average that builds the detached error maps for UOT.

During mixed training, the cached ROI can be randomly shifted by at most $3\%$
of its own width/height and enlarged by a factor sampled from $[1,1.15]$.  Pass~A
and Pass~R then receive independent random affine and intensity transforms.  The
default affine ranges are $\pm10^\circ$ rotation, scale $[0.95,1.05]$,
translation $\pm5\%$ of the image extent, horizontal-flip probability $0.5$, and
vertical-flip probability $0.2$.  Intensity perturbation is applied with
probability $0.8$: contrast is sampled from $[0.8,1.2]$, the intensity shift from
$[-0.1,0.1]$ times the slice standard deviation, and Gaussian-noise scale from
$[0,0.04]$ times that standard deviation.  The BiomedParse image additionally
uses a gamma sampled from $[0.75,1.35]$ after clipping to $[0,255]$.

\subsection{Two-pass loss}

During the anatomy warm-up, the LGE branch objective is
\begin{equation}
    \cL_{A}=\cL^{A}_{\mathrm{seg}}
      +0.2\cL^{A}_{\mathrm{ae}}
      +0.3\cL^{A}_{\mathrm{ort}}
      +\lambda_{\mathrm{route}}(e)\cL^{A}_{\mathrm{route}}
      +\lambda_p(e)\cL^{A}_{P\text{-OT}}
      +\lambda_s(e)\cL^{A}_{S\text{-UOT}}.
\end{equation}
Once ROI training starts, each branch's non-segmentation terms are multiplied by
$1/2$ because they are computed twice.  With the implemented
$\lambda_A=\lambda_R=1$, the joint objective is
\begin{equation}
\begin{split}
    \cL_{\mathrm{LGE}}={}&\lambda_A\left[
      \cL^{A}_{\mathrm{seg}}+\tfrac12\cL^{A}_{\mathrm{aux}}
      \right]
      +\lambda_R\left[
      \cL^{R}_{\mathrm{seg}}+\tfrac12\cL^{R}_{\mathrm{aux}}
      \right],\\
    \cL^{j}_{\mathrm{aux}}={}&
      0.2\cL^{j}_{\mathrm{ae}}
      +0.3\cL^{j}_{\mathrm{ort}}
      +\lambda_{\mathrm{route}}(e)\cL^{j}_{\mathrm{route}}
      +\lambda_p(e)\cL^{j}_{P\text{-OT}}
      +\lambda_s(e)\cL^{j}_{S\text{-UOT}},
      \quad j\in\{A,R\}.
    \label{eq:lge-joint-loss}
\end{split}
\end{equation}
The two segmentation terms are not halved.  Both forwards are accumulated before
one optimizer update.

\subsection{Spatial hard switching at inference}

The refinement logits are bilinearly resized from the model-size ROI back to the
ROI's full-image extent.  Let $r(\bm{u})$ be its binary rectangular indicator.
For the deployed V2/V3 decision rule, foreground logit $c$ is
\begin{equation}
    \ell^{\mathrm{final}}_c(\bm{u})=
    \begin{cases}
      \ell^{R}_c(\bm{u}),&r(\bm{u})=1,\\
      \ell^{A}_c(\bm{u}),&r(\bm{u})=0,
          \;c\in\{\mathrm{LV},\mathrm{RV}\},\\
      -20,&r(\bm{u})=0,
          \;c\notin\{\mathrm{LV},\mathrm{RV}\}.
    \end{cases}
    \label{eq:hard-switch}
\end{equation}
There is no cross-pass averaging or soft blending.  A zero background logit is
prepended.  Let $\mathcal{R}_{\mathrm{logit}}$ denote the composite logit-space
restoration: in-plane bilinear resizing to the aligned preprocessed geometry,
followed by the nnUNet probability-resampling function to the
pre-cropping geometry.  The exclusive prediction in that geometry is
\begin{equation}
    \widehat Y_{\mathrm{crop}}(\bm{u})
    =\arg\max_{c\in\{0,1,\ldots,C_R\}}
       \mathcal{R}_{\mathrm{logit}}\!\left(
       0,\ell^{\mathrm{final}}_1,\ldots,
       \ell^{\mathrm{final}}_{C_R}\right)_c(\bm{u}).
    \label{eq:exclusive-prediction}
\end{equation}
Finally, the discrete segmentation is inserted into the stored crop bounding box
and inverse-transposed to raw NIfTI geometry.  Thus, the spatial hard switch is
formed on the $256\times256$ model grid, but class argmax occurs only after logit
resampling to the original cropped geometry.

\section{Training--Deployment Boundary}

% Code anchors: oodka/train/forward.py::predict_block_logits_per_class;
%               oodka/train/engine.py::_save_checkpoint; run_eval_oodka.py;
%               run_eval_lge_roi.py

The expert is strictly training-only.  At deployment, the system loads
\begin{equation}
    \{B_{\phi_b},D^b_{\ell,p},D^b_{\ell,s}\}_{\ell\in\mathcal{I}}
    \quad\text{and the spatial Beta router},
\end{equation}
but does not load nnUNet, the expert autoencoders, any expert preprocessing
tensor, or any OT/UOT module.  BiomedParse features are decomposed, decoded,
re-fused using the deterministic Beta mean in Eq.~\eqref{eq:beta-mean}, and sent
to the frozen predictor.  LGE deployment runs this pure-student forward twice,
first on the full slice and then on its predicted ROI.

\section{Implementation Constants and Experiment Profile}
For the current $256\times256$ LGE instantiation, the four spatial feature sizes
are $(64,32,16,8)^2$.  The nnUNet channel counts are
$(128,256,512,512)$ and the BiomedParse channel counts are
$(192,384,768,1536)$ for levels $2$--$5$, respectively.  Consequently, the OT
grids are $(32,32,16,8)^2$ after the per-dimension cap.  The trainable modules
contain $10{,}140{,}994$ parameters in total; the deployment subset (four
foundation decomposition modules plus the router) contains $6{,}432{,}130$.
\end{document}
